\documentclass[12pt,letterpaper]{article}
\usepackage[utf8]{inputenc}
\usepackage[spanish]{babel}
\usepackage{amsmath}
\usepackage{amsfonts}
\usepackage{amssymb}
\usepackage[margin=2.5cm]{geometry}

\begin{document}

\begin{center}
    \Large \textbf{Evaluación de Examen 2: Unidad II} \\
    \large Física para Ciencias de la Salud (005-1713) \\
    \vspace{0.5cm}
    \textbf{Estudiante:} Rosmeilys Sucre \quad \textbf{C.I.:} 33.386.357
    \vspace{0.3cm} \\
    \large \textbf{Calificación Final: 9/10 puntos}
\end{center}

\vspace{0.5cm}

\section*{Análisis de Resultados}

\subsection*{1. Cinemática (Aceleración media)}
\textbf{Procedimiento y resultado:} Plantea correctamente la ecuación de aceleración media restando las velocidades y dividiendo entre el tiempo ($0,6 \text{ m/s} / 0,15 \text{ s}$). El cálculo aritmético es exacto ($4$); sin embargo, falla en el análisis dimensional al encuadrar el resultado final con unidades de velocidad ($4 \text{ m/s}$) en lugar de aceleración ($4 \text{ m/s}^2$). \\
\textbf{Veredicto:} \textbf{Parcialmente correcto} (Penalización leve por error en las unidades de la respuesta final) (1,5/2 pts).

\subsection*{2. Dinámica (Leyes de Newton)}
\textbf{Procedimiento y resultado:} Muestra un planteamiento impecable partiendo de la Segunda Ley de Newton para hallar la aceleración ($a = \frac{F}{m}$). Sustituye los valores dividiendo $150 \text{ N}$ entre $25 \text{ kg}$, logrando un resultado correcto de $6 \text{ m/s}^2$. No obstante, al transcribir su respuesta final al lado derecho de la hoja, anota erróneamente $6 \text{ m/s}$, reiterando la confusión con las unidades. \\
\textbf{Veredicto:} \textbf{Parcialmente correcto} (Penalización por transcripción errónea de la unidad final) (1,5/2 pts).

\subsection*{3. Trabajo Mecánico}
\textbf{Procedimiento y resultado:} Identifica correctamente la fórmula general del trabajo, incluyendo el componente trigonométrico ($W = F \cdot d \cdot \cos(0^\circ)$). Efectúa la multiplicación de los factores ($400 \text{ N} \cdot 0,8 \text{ m} \cdot 1$) obteniendo de forma perfectamente exacta $320 \text{ J}$. \\
\textbf{Veredicto:} \textbf{Correcto} (2/2 pts).

\subsection*{4. Energía Potencial}
\textbf{Procedimiento y resultado:} Extrae los datos (masa, gravedad y altura) y plantea el producto de las tres variables según la ecuación de la energía potencial gravitatoria ($E_p = 1,2 \text{ kg} \cdot 9,8 \text{ m/s}^2 \cdot 1,5 \text{ m}$). Efectúa el cálculo matemático obteniendo el resultado verificado y correcto de $17,64 \text{ J}$. \\
\textbf{Veredicto:} \textbf{Correcto} (2/2 pts).

\subsection*{5. Potencia Mecánica}
\textbf{Procedimiento y resultado:} Relaciona correctamente el trabajo efectuado y el tiempo mediante la fracción respectiva ($P = \frac{75 \text{ J}}{60 \text{ s}}$). El desarrollo de la división arroja la cifra correcta y verificada de $1,25 \text{ Watts}$. \\
\textbf{Veredicto:} \textbf{Correcto} (2/2 pts).

\vspace{0.5cm}
\section*{Comentarios Generales y Sugerencias}
¡Un excelente examen! La estudiante demuestra un dominio muy sólido de los conceptos de la Unidad II, logrando aplicar las fórmulas correctas y ejecutar cálculos matemáticos impecables en cada uno de los problemas.
\begin{itemize}
    \item Existe una excelente claridad en los desarrollos aritméticos, así como un estupendo entendimiento para los bloques de Trabajo, Energía y Potencia, donde se manejaron a la perfección los Joules y Watts.
    \item Se aconseja tener mucha más precaución con las \textbf{unidades del Sistema Internacional} en la cinemática. En física, la aceleración siempre debe llevar un exponente cuadrado en la variable del tiempo ($\text{m/s}^2$). Omitirlo altera por completo la magnitud clínica y científica del resultado.
    \item Como un pequeño detalle ortográfico en la notación técnica: el símbolo de \textit{Watts} se puede escribir como ``W'' (en mayúscula), o bien, si se escribe la palabra completa en español, lo ideal es mantener el singular.
    \item ¡Sigue con este mismo entusiasmo y excelente nivel analítico!
\end{itemize}

\end{document}